%% kapa: exam=2 round=32 year=2021 period=1 subject=appraisal_practice subject_ko=감정평가실무 seq=3 number=3 points=20
%% source: Q-Net 감정평가사 기출문제 (cst00309) artlSeq=5208993 file=raw/1008/5208993/감평 2차 시험지(2021).pdf
%% conversion: kordoc 4.24.51 -> tools/gen2.py (auto); review=proofread-sonnet; flags=none
\begin{longquestion}{【문제 3】}
감정평가사 甲은 중앙토지수용위원회로부터 이의재결평가를 의뢰받았다. 관련 법규 및 이론을 참작하고 제시된 자료를 활용하여 적정보상액을 산정하시오. (20점) \par
\medskip\noindent < 자료 1 > 사업개요 \par
\medskip\noindent 1. 사업명: ○○민자고속화도로사업 \par
2. 사업시행자: ○○민자고속화도로 주식회사 \par
3. 사업인정고시(의제)일: 2020년 10월 2일 \par
4. 수용재결일: 2021년 4월 1일 \par
5. 의뢰인: 2021년 8월 6일 \par
\medskip\noindent < 자료 2 > 의뢰목록 \par
\medskip\noindent 1. 토지 목록 \par
{\small\centering
\begin{tabular}{|c|c|c|c|c|c|c|c|}\hline
\makecell{일련\\번호} & 소재지 & 지번 & 지목 & \makecell{이용\\상황} & 용도지역 & \makecell{편입면적\\(m\textsuperscript{2})} & 피수용자 \\\hline
1 & \makecell{A군\\B읍 C리} & 산1 & 임야 & 전 & 자연녹지지역 & 3,000 & A군 \\\hline
2 & \makecell{A군\\B읍 C리} & 산2 & 임야 & 자연림 & \makecell{자연녹지지역\\보전녹지지역} & 5,000×1/2 & 乙 \\\hline
\end{tabular}\par}
2. 지장물 목록 \par
{\small\centering
\begin{tabular}{|c|c|c|c|c|c|c|}\hline
\makecell{일련\\번호} & 소재지 & 지번 & 물건종류 & 규격 & 수량 & 피수용자 \\\hline
3 & \makecell{A군\\B읍 C리} & 산1 & 개간비 & 개간면적(3,000m\textsuperscript{2}) & 1식 & 丙 \\\hline
\end{tabular}\par}
\medskip\noindent < 자료 3 > 인근지역 공시지가표준지(공시기준일: 2020년 1월 1일) \par
{\small\centering
\resizebox{\linewidth}{!}{\begin{tabular}{|c|c|c|c|c|c|c|c|c|c|}\hline
\makecell{기\\호} & 소재지 & 지번 & \makecell{면적\\(m\textsuperscript{2})} & 지목 & \makecell{이용\\상황} & \makecell{용도\\지역} & \makecell{도로\\접면} & \makecell{형상\\지세} & \makecell{공시지가\\(원/m\textsuperscript{2})} \\\hline
A & \makecell{A군\\B읍 C리} & 10 & 1,820 & 전 & 대 & \makecell{자연녹지\\지역} & \makecell{세로\\(가)} & \makecell{세장형\\완경사} & 350,000 \\\hline
B & \makecell{A군\\B읍 C리} & 20 & 950 & 답 & 전 & \makecell{자연녹지\\지역} & \makecell{세로\\(불)} & \makecell{부정형\\완경사} & 120,000 \\\hline
C & \makecell{A군\\B읍 C리} & 45 & 8,452 & 임야 & 조림 & \makecell{자연녹지\\지역} & \makecell{세로\\(불)} & \makecell{부정형\\완경사} & 60,000 \\\hline
D & \makecell{A군\\B읍 C리} & 산10-1 & 5,526 & 임야 & \makecell{자연\\림} & \makecell{보전녹지\\지역} & 맹지 & \makecell{부정형\\완경사} & 35,000 \\\hline
E & \makecell{A군\\B읍 C리} & 산15 & 2,570 & 임야 & \makecell{자연\\림} & \makecell{자연녹지\\지역} & 맹지 & \makecell{부정형\\급경사} & 15,000 \\\hline
\end{tabular}}\par}
\medskip\noindent < 자료 4 > 인근 평가사례 및 매매사례 \par
\medskip\noindent 1. 평가사례(ㄱ) \par
- 지목, 이용상황, 면적: 전, 전, 2,570m\textsuperscript{2} \par
- 용도지역: 자연녹지지역 \par
- 평가목적: 담보 \par
- 기준시점: 2020년 8월 1일 \par
- 평가단가: 220,000원/m\textsuperscript{2} \par
- 기타사항: 임야지대 내 적법하게 개간된 전으로 이용 중인 사례임 \par
\medskip\noindent 2. 평가사례(ㄴ) \par
- 지목, 이용상황, 면적: 전, 답, 416m\textsuperscript{2} \par
- 용도지역: 자연녹지지역 \par
- 평가목적: 협의보상 \par
- 기준시점: 2020년 12월 1일 \par
- 평가단가: 270,000원/m\textsuperscript{2} \par
- 기타사항: 대상 공익사업에 포함된 협의완료 된 사례임 \par
\medskip\noindent 3. 평가사례(ㄷ) \par
- 지목, 이용상황, 면적: 임야, 자연림, 5,470m\textsuperscript{2} \par
- 용도지역: 자연녹지지역 \par
- 평가목적: 체납처분 \par
- 기준시점: 2021년 3월 1일 \par
- 평가단가: 130,000원/m\textsuperscript{2} \par
- 기타사항: 유찰사례로 처분절차 진행중인 사례임 \par
\medskip\noindent 4.평가사례(ㄹ) \par
- 지목, 이용상황, 면적: 임야, 자연림, 1,320m\textsuperscript{2} \par
- 용도지역: 보전녹지지역 \par
- 평가목적: 협의보상 \par
- 기준시점: 2020년 9월 1일 \par
- 평가단가: 75,000원/m\textsuperscript{2} \par
- 기타사항: 임지상에 소재하는 잡목을 포함한 일괄 평가사례임 \par
\medskip\noindent 5. 거래사례(ㅁ) \par
- 지목, 이용상황, 면적: 전, 전, 1,560m\textsuperscript{2} \par
- 용도지역: 자연녹지지역 \par
- 거래시점: 2020년 7월 31일 \par
- 총 거래가격: 399,360,000원 \par
- 기타사항: 개인과 법인간의 거래사례임 \par
\medskip\noindent 6. 거래사례(ㅂ) \par
- 지목, 이용상황 면적: 전, 전, 1,906m\textsuperscript{2} \par
- 용도지역: 자연녹지지역 \par
- 거래시점: 2020년 2월 1일 \par
- 총 거래가격: 590,860,000원 \par
- 기타사항: 친족 간의 지분거래사례임 \par
\medskip\noindent 7. 거래사례(ㅅ) \par
- 지목, 이용상황 면적: 임야, 자연림, 3,750m\textsuperscript{2} \par
- 용도지역: 자연녹지지역 \par
- 거래시점: 2020년 7월 1일 \par
- 총 거래가격: 562,500,000원 \par
- 기타사항: 임지상에 소재하는 잣나무(300그루)를 포함한 일괄 거래사례임 \par
\medskip\noindent 8. 거래사례(ㅇ) \par
- 지목, 이용상황 면적: 임야, 자연림, 1,670m\textsuperscript{2} \par
- 용도지역: 보전녹지지역 \par
- 거래시점: 2021년 3월 1일 \par
- 총 거래가격: 158,650,000원 \par
- 기타사항: 최근 지가상승이 반영된 정상거래사례임 \par
\medskip\noindent < 자료 5 > 지가변동률(A군, 녹지지역) \par
{\small\centering
\resizebox{\linewidth}{!}{\begin{tabular}{|c|c|c|c|c|c|c|c|c|c|c|c|c|c|}\hline
2020년 & 1월 & 2월 & 3월 & 4월 & 5월 & 6월 & 7월 & 8월 & 9월 & 10월 & 11월 & 12월 & 누계액 \\\hline
변동률 & 0.092 & 0.313 & 0.223 & 0.252 & 0.252 & 0.170 & 0.363 & 0.230 & 0.280 & 0.223 & 0.223 & 0.312 & 2.972 \\\hline
\end{tabular}}\par}
\vspace{0.5em}
\par{\small\centering
\resizebox{\linewidth}{!}{\begin{tabular}{|c|c|c|c|c|c|c|c|c|c|c|c|c|}\hline
2021년 & 1월 & 2월 & 3월 & 4월 & 5월 & 6월 & 7월 & 8월 & 9월 & 10월 & 11월 & 12월 \\\hline
변동률 & 0.282 & 0.221 & 0.235 & 0.310 & 0.289 & 0.287 & 미고시 & 미고시 & 미고시 & 미고시 & 미고시 & 미고시 \\\hline
\end{tabular}}\par}
\medskip\noindent < 자료 6 > 요인 격차율 \par
\medskip\noindent 1. 지역요인: 대상과 사례는 인근지역에 소재하여 지역요인은 유사함 \par
2. 개별요인 \par
{\small\centering
\begin{tabular}{|c|c|c|}\hline
구분 & 일련번호 1 & 일련번호 2 \\\hline
공시지가표준지 A & 0.90 & 0.65 \\\hline
공시지가표준지 B & 1.05 & 0.75 \\\hline
공시지가표준지 C & 1.02 & 1.03 \\\hline
공시지가표준지 D & 1.50 & 1.08 \\\hline
공시지가표준지 E & 1.15 & 1.10 \\\hline
\end{tabular}\par}
\vspace{0.5em}
\par{\small\centering
\begin{tabular}{|c|c|c|c|c|c|}\hline
구분 & \makecell{공시지가\\표준지 A} & \makecell{공시지가\\표준지 B} & \makecell{공시지가\\표준지 C} & \makecell{공시지가\\표준지 D} & \makecell{공시지가\\표준지 E} \\\hline
평가사례(ㄱ) & 0.85 & 0.90 & 0.55 & 0.35 & 0.40 \\\hline
평가사례(ㄴ) & 0.90 & 0.95 & 0.60 & 0.40 & 0.45 \\\hline
평가사례(ㄷ) & 0.65 & 0.70 & 0.95 & 0.65 & 0.70 \\\hline
평가사례(ㄹ) & 0.50 & 0.55 & 1.30 & 0.90 & 0.80 \\\hline
거래사례(ㅁ) & 0.88 & 0.85 & 0.58 & 0.38 & 0.50 \\\hline
거래사례(ㅂ) & 0.80 & 0.78 & 0.50 & 0.30 & 0.43 \\\hline
거래사례(ㅅ) & 0.67 & 0.75 & 1.25 & 0.60 & 0.65 \\\hline
거래사례(ㅇ) & 0.55 & 0.60 & 1.15 & 0.88 & 0.68 \\\hline
\end{tabular}\par}
\medskip\noindent < 자료 7 > 기타자료 \par
\medskip\noindent 1. 대상은 20만m\textsuperscript{2} 미만 공익사업으로 협의와 수용재결 절차가 완료된 상태로 일부 피수용자에 대한 이의재결이 진행중임 \par
2. 일련번호(1), (2)의 지세는 완경사지임 \par
2. 丙은 일련번호(1)을 2018년 10월 2일부터 관계법령에 따라 적법하게 개간하고 현재까지 적법하게 점유하고 있으며, 개간소요비용은 개간당시 250,000,000원이나, 가격시점 기준 300,000,000원이 소요됨 % [원문 그대로: 번호 2 중복]
\par
3. 일련번호(2)의 자연녹지지역 비율은 전체의 60\%임 \par
4. 인근지역은 잣나무만의 거래가 일반적이며, 1그루당 500,000원에 거래됨 \par
5. 토지단가 산출시 백 원 단위에서 반올림 할 것 \par
\end{longquestion}
